% ============================================================
\section{Contexto: el grupo fundamental}
% ============================================================

\begin{frame}{¿Qué estamos construyendo?}
  \begin{block}{Definición — $\pi_1(X,\xo)$}
    El \textbf{grupo fundamental} de $X$ en $\xo$ es el conjunto de clases de
    homotopía de lazos con base $\xo$, con la operación
    \[
      \clase{\alpha} \circ \clase{\beta} \;=\; \clase{\alpha * \beta}
    \]
    donde $(\alpha * \beta)(t) =
      \begin{cases} \alpha(2t) & 0 \le t \le \tfrac{1}{2}\\
                    \beta(2t-1) & \tfrac{1}{2} \le t \le 1 \end{cases}$
  \end{block}

  \bigskip
  \textbf{Teorema 4.2:} $\pione(X,\xo)$ es un \emph{grupo} bajo $\circ$.
  \medskip

  \begin{columns}[T]
    \begin{column}{0.3\textwidth}
      \centering
      \textcolor{verdeTeal}{\textbf{Lema A}}\\\smallskip
      Identidad $\clase{c}$\\\smallskip
      {\footnotesize (ya demostrado)}
    \end{column}
    \begin{column}{0.3\textwidth}
      \centering
      \textcolor{morado}{\textbf{Lema B}}\\\smallskip
      Inversos $\clase{\abar}$\\\smallskip
      {\footnotesize \textbf{← hoy}}
    \end{column}
    \begin{column}{0.3\textwidth}
      \centering
      \textcolor{coralOscuro}{\textbf{Lema C}}\\\smallskip
      Asociatividad\\\smallskip
      {\footnotesize \textbf{← hoy}}
    \end{column}
  \end{columns}
\end{frame}

\begin{frame}{Herramienta clave: Lema de Continuidad}
  \begin{alertblock}{Lema de Continuidad}
    Sea $X = A \cup B$ con $A, B$ cerrados. Si $f:A \to Y$ y $g:B \to Y$ son
    continuas con $f(x)=g(x)$ para todo $x \in A \cap B$, entonces
    \[
      h(x) = \begin{cases} f(x) & x \in A \\ g(x) & x \in B \end{cases}
    \]
    es continua.
  \end{alertblock}

  \bigskip
  \textbf{Uso práctico en este capítulo:} todas las homotopías se definen por
  tramos sobre el cuadrado $\IxI$.  El Lema garantiza la continuidad cuando
  las dos piezas coinciden en la frontera común.
\end{frame}
